\subsection{RQ6: How does an iterative static-analysis feedback cycle affect the ability of language models to generate correct code?}
\label{sec:rq6}

This RQ compares the final code produced by the iterative tool-assisted protocol with the original generation results of RQ1. The comparison concerns two complete generation protocols. RQ1 uses a temperature of 1.0 and a single generation, whereas the iterative protocol uses a temperature of 0.0 together with static-analysis feedback, self-review, and up to five revision cycles. Consequently, the observed differences cannot be attributed exclusively to static analysis and are interpreted as changes under the complete iterative protocol.

\begin{table}[htbp]
\centering
\caption{Code generation under the iterative static-analysis protocol, compared with RQ1. Values are percentages and the second line in each outcome cell reports the change relative to RQ1. Green changes indicate an improvement in the corresponding outcome, while red changes indicate a deterioration. Avg Cyc is the mean number of cycles.}
\label{tab:rq6_generation_delta}
\setlength{\tabcolsep}{1pt}

\newcommand{\better}[1]{\textcolor{green!55!black}{#1}}
\newcommand{\worse}[1]{\textcolor{red!70!black}{#1}}
\newcommand{\lightmodelrule}{%
  \arrayrulecolor{black!15}\specialrule{0.25pt}{0.5pt}{0.5pt}\arrayrulecolor{black}%
}

\begin{tabular}{@{}>{\centering\arraybackslash}m{0.29\columnwidth}*{5}{>{\centering\arraybackslash}m{0.132\columnwidth}}@{}}
\toprule
\small\textbf{Model} &
\scriptsize\textbf{Avg Cyc} &
\scriptsize\textbf{Correct} &
\scriptsize\textbf{Funct.} &
\scriptsize\textbf{Runtime} &
\scriptsize\textbf{Syntax} \\
\midrule
\multicolumn{6}{l}{\emph{Small language models}} \\
\midrule
\shortstack[c]{Gemma 4\\(31B)}
& 1.29
& \shortstack{82\\\better{+3}}
& \shortstack{15\\0}
& \shortstack{4\\0}
& \shortstack{0\\\better{-3}} \\
\lightmodelrule
\shortstack[c]{GPT-OSS\\(20B)}
& 1.46
& \shortstack{80\\\better{+1}}
& \shortstack{15\\\better{-2}}
& \shortstack{4\\\better{-1}}
& \shortstack{2\\\worse{+1}} \\
\lightmodelrule
\shortstack[c]{Muse Glimmer\\(30B)}
& 1.41
& \shortstack{79\\\worse{-2}}
& \shortstack{16\\\worse{+1}}
& \shortstack{4\\0}
& \shortstack{2\\\worse{+2}} \\
\lightmodelrule
\shortstack[c]{Qwen 3.8\\(27B)}
& 1.17
& \shortstack{78\\\better{+2}}
& \shortstack{19\\\better{-1}}
& \shortstack{4\\\better{-1}}
& \shortstack{0\\0} \\
\lightmodelrule
\shortstack[c]{Qwen 2.5 Coder\\(32B)}
& 1.36
& \shortstack{75\\\better{+8}}
& \shortstack{20\\\better{-5}}
& \shortstack{7\\\better{-2}}
& \shortstack{0\\0} \\
\lightmodelrule
\shortstack[c]{Qwen 3 Coder\\(30B)}
& 1.57
& \shortstack{74\\\worse{-2}}
& \shortstack{21\\\worse{+4}}
& \shortstack{6\\\better{-1}}
& \shortstack{0\\0} \\
\lightmodelrule
\shortstack[c]{Devstral Small 2\\(24B)}
& 1.37
& \shortstack{73\\\better{+2}}
& \shortstack{22\\\better{-2}}
& \shortstack{6\\\better{-1}}
& \shortstack{0\\0} \\
\lightmodelrule
\shortstack[c]{Ornith 1.5\\(35B)}
& 1.68
& \shortstack{72\\\better{+6}}
& \shortstack{19\\\worse{+1}}
& \shortstack{6\\\better{-1}}
& \shortstack{5\\\better{-6}} \\
\lightmodelrule
\shortstack[c]{Qwen 2.5 Coder\\(7B)}
& 2.06
& \shortstack{71\\\worse{-3}}
& \shortstack{24\\\worse{+3}}
& \shortstack{8\\0}
& \shortstack{0\\0} \\
\lightmodelrule
\shortstack[c]{GLM 4.7 Flash\\(30B)}
& 2.14
& \shortstack{63\\\worse{-3}}
& \shortstack{25\\\better{-4}}
& \shortstack{5\\\better{-2}}
& \shortstack{8\\\worse{+8}} \\
\lightmodelrule
\shortstack[c]{Gemma 3\\(4B)}
& 1.43
& \shortstack{61\\\worse{-2}}
& \shortstack{34\\\worse{+1}}
& \shortstack{8\\\better{-1}}
& \shortstack{0\\0} \\
\midrule
\multicolumn{6}{l}{\emph{Large language models}} \\
\midrule
\shortstack[c]{DeepSeek V4.1\\Flash\\(552B)}
& 1.14
& \shortstack{81\\\better{+7}}
& \shortstack{16\\\better{-1}}
& \shortstack{5\\\better{-1}}
& \shortstack{0\\\better{-6}} \\
\lightmodelrule
\shortstack[c]{MiMo V2.5\\(310B)}
& 1.15
& \shortstack{80\\\better{+7}}
& \shortstack{16\\\better{-2}}
& \shortstack{5\\\better{-2}}
& \shortstack{0\\\better{-2}} \\
\bottomrule
\end{tabular}
\end{table}

The effect of the iterative protocol on final correctness is heterogeneous among the SLMs. Six of the eleven evaluated SLMs increase their proportion of \emph{Correct} solutions relative to RQ1, while five decrease it. The largest positive change is observed for Qwen 2.5 Coder 32B, which increases Correct by approximately 8 percentage points, followed by Ornith 1.5 with approximately 6 points and Gemma 4 31B with approximately 3 points. In contrast, Qwen 2.5 Coder 7B, GLM 4.7 Flash, Qwen 3 Coder, Muse Glimmer, and Gemma 3 finish the iterative protocol with fewer Correct solutions than under the original generation condition. Static-analysis-assisted revision therefore does not produce a uniform improvement in final code correctness across SLMs.

The error distribution shows that different models reach their final Correct proportion through different changes in failure type. Qwen 2.5 Coder 32B provides the clearest positive example. Its increase of approximately 8 points in Correct is accompanied by reductions of approximately 5 points in \emph{Functional Error} and 2 points in \emph{Runtime Error}, while Syntax remains near zero. The largest SLM improvement therefore cannot be explained only by correcting properties directly checked by compilation or Pyright. The iterative protocol can also change Functional outcomes even though the static tools provide no direct functional-correctness signal.

Ornith 1.5 presents a different improvement profile. Its Correct proportion increases by approximately 6 points, while Syntax decreases by approximately 6 points and Runtime decreases slightly. Functional Error does not improve and instead increases slightly. Gemma 4 31B shows an even more direct relationship between structural failures and final correctness: Correct increases by approximately 3 points while Syntax decreases by approximately the same amount, with comparatively small changes in the remaining outcomes. Thus, similar gains in Correct may originate from different changes in the underlying error distribution.

The GLM 4.7 Flash result illustrates the opposite phenomenon. Functional and Runtime failures decrease by approximately 4 and 2 percentage points, respectively, yet Syntax increases by approximately 8 points and the final proportion of Correct solutions decreases by approximately 3 points. Improving some observable failure categories therefore does not guarantee a better final result if the protocol simultaneously introduces another type of failure. This case is particularly notable because GLM reaches 94.1\% Syntax Recall as a detector in RQ2, but produces substantially more Syntax failures under the iterative generation protocol. Recognizing a syntactic problem during classification and avoiding or repairing the same type of problem during generation are therefore distinct capabilities.

\emph{Runtime Error} presents an important contrast with RQ5. When static feedback is supplied for classification, Runtime Recall decreases for most detector configurations. During iterative generation, however, the incidence of Runtime failures decreases for eight of the eleven SLMs and remains approximately unchanged for the other three in the rounded results reported in Table~\ref{tab:rq6_generation_delta}. The two substantially larger models also reduce Runtime failures. Qwen 2.5 Coder 32B presents the largest SLM reduction, followed by smaller decreases for GLM 4.7 Flash, Qwen 3.8, Ornith 1.5, GPT-OSS, Devstral Small 2, Qwen 3 Coder, and Gemma 3. This difference suggests that translating partial static evidence into an explicit \emph{Runtime Error} label and using similar evidence as guidance for code revision are different tasks.

The protocol also affects Functional outcomes in both directions. Qwen 2.5 Coder 32B, GLM 4.7 Flash, GPT-OSS, Qwen 3.8, and Devstral Small 2 reduce Functional failures, while Qwen 3 Coder, Qwen 2.5 Coder 7B, Muse Glimmer, Ornith 1.5, and Gemma 3 increase them. This behavior is important because neither compilation nor Pyright directly evaluates whether the candidate satisfies the functional specification. Functional changes can therefore emerge from the model's self-review and revision behavior even when the external evidence concerns only syntactic or statically detectable runtime properties.

The average number of cycles does not show a simple relationship with the final outcome. Although the protocol permits up to five rounds, the mean number of cycles ranges only from 1.14 to 2.14 across the evaluated models. Qwen 2.5 Coder 32B obtains the largest SLM increase in Correct with an average of only 1.36 cycles, whereas GLM 4.7 Flash averages 2.14 cycles and finishes with fewer Correct solutions. Qwen 2.5 Coder 7B similarly averages 2.06 cycles while decreasing Correct. Additional revision opportunities therefore do not by themselves imply better final code; the result depends on how effectively each model uses the available feedback during revision.

For the two hosted models, paired observations are available for both the original and iterative protocols. DeepSeek changes from 13 tasks that are Correct only under RQ1 to 54 tasks that are Correct only under the iterative protocol, while MiMo changes from 11 RQ1-only Correct tasks to 48 iterative-only Correct tasks. Exact two-sided McNemar tests yield $p<.001$ for both models. The retained SLM comparisons remain descriptive because equivalent paired raw observations are not available for those conditions.

Overall, the iterative protocol improves final correctness for some models but produces regressions for others. Its effects are also not restricted to the properties directly inspected by the static tools: some of the largest gains and losses involve Functional Error, while Runtime failures decrease more consistently during revision than Runtime Recall improves during classification in RQ5. The results therefore reinforce the distinction between error detection, interpretation of external evidence, and code revision. Because RQ1 and RQ6 also differ in decoding temperature, self-review, and the opportunity for repeated revision, these results characterize the complete iterative protocol rather than the isolated causal contribution of static-analysis feedback.